В данной лабораторной работе были использованы следующие технологии:

\begin{itemize}
    \item \textbf{Python} -- язык программирования высокого уровня, используемый для создания backend-сервера, обработки данных и реализации алгоритмов машинного обучения для прогнозирования водных ресурсов.
    
    \item \textbf{FastAPI} -- современный веб-фреймворк для создания REST API на Python, обеспечивающий высокую производительность и автоматическую генерацию документации.
    
    \item \textbf{Uvicorn} -- ASGI-сервер для запуска FastAPI приложений, обеспечивающий асинхронную обработку запросов.
    
    \item \textbf{Pandas} -- библиотека Python для анализа и обработки данных, используемая для нормализации Excel-файлов в единый CSV-формат и работы с временными рядами.
    
    \item \textbf{NumPy} -- библиотека Python для численных вычислений, необходимая для математических операций при обработке данных и построении прогнозов.
    
    \item \textbf{scikit-learn} -- библиотека машинного обучения для Python, используемая для реализации алгоритмов прогнозирования: полиномиальная регрессия, Ridge-регрессия, Kernel Ridge и сплайны.
    
    \item \textbf{Prophet} -- библиотека от Facebook для прогнозирования временных рядов, используемая как альтернативный метод построения прогнозов с учетом сезонности.
    
    \item \textbf{TypeScript} -- типизированный надмножество JavaScript, используемое для разработки frontend-приложения с улучшенной надежностью кода.
    
    \item \textbf{Next.js} -- React-фреймворк для создания полнофункциональных веб-приложений с серверным рендерингом и оптимизацией производительности.
    
    \item \textbf{React} -- библиотека JavaScript для построения пользовательских интерфейсов, используемая для создания интерактивных компонентов веб-приложения.
    
    \item \textbf{Leaflet} -- открытая JavaScript-библиотека для создания интерактивных карт, используемая для визуализации географических данных о реках и водных объектах Беларуси.
    
    \item \textbf{react-leaflet} -- React-обертка для библиотеки Leaflet, обеспечивающая интеграцию картографических компонентов в React-приложение.
    
    \item \textbf{Recharts} -- библиотека React для создания диаграмм и графиков, используемая для визуализации временных рядов и прогнозных данных водных ресурсов.
    
    \item \textbf{Axios} -- HTTP-клиент для JavaScript, используемый для выполнения запросов к backend API и получения данных о водных ресурсах.
    
    \item \textbf{SWR} -- библиотека React Hooks для получения данных, обеспечивающая кэширование, ревалидацию и автоматическое обновление данных.
    
    \item \textbf{Tailwind CSS} -- utility-first CSS-фреймворк для быстрого создания современного дизайна интерфейса с эффектом стеклянного морфизма (glassmorphism).
    
    \item \textbf{GeoJSON} -- формат данных для представления географических объектов, используемый для хранения и передачи контуров рек и водных бассейнов.
    
    \item \textbf{Excel (xlsx)} -- формат файлов Microsoft Excel, используемый как исходный формат данных о водных ресурсах, которые затем нормализуются в CSV.
    
    \item \textbf{CSV} -- текстовый формат для представления табличных данных, используемый как унифицированный формат хранения нормализованных данных после обработки Excel-файлов.
\end{itemize}

